% sec-03: cost efficiency and clinical evidence, small-business framing.
\section{What the Company Has Already Demonstrated}

The strongest commercial evidence a pre-award small business can offer is its
own unit economics. Industry benchmarks put a single virtual trial run above
\$120,000 and 4.5 months, and a quantitative systems pharmacology virtual trial
above \$2 million and several months to a year or more. The three simulations
here cost an estimated \$36,330 each and took about a month each
\cite{daraxsims}.

\begin{apptable}
\begin{tabularx}{\textwidth}{>{\raggedright\arraybackslash}p{3.9cm} >{\raggedright\arraybackslash}p{3.5cm} >{\raggedright\arraybackslash}X}
\headrow \thead{Work product} & \thead{Industry benchmark} & \thead{This company}\\
\midrule
Empirical virtual trial, triplicate & Above \$120,000, 4.5 months per run & \$36,330 estimated, about one month \\
Ten-arm QSP virtual trial & Above \$2,000,000, months to a year & \$36,330 estimated, about one month \\
Digital-twin simulator platform & \$28,000 to \$700,000, 3 to 16 weeks & \$36,330 estimated, about one month \\
\end{tabularx}
\end{apptable}
\tabcap{Three work products against their industry cost and schedule
benchmarks. The comparison is the company's demonstrated burn rate, which is
what a Phase I reviewer is actually being asked to underwrite.}

On the clinical side, the August 2025 quantitative systems pharmacology run
returned median overall survival of \textbf{12.8 months} for the daraxonrasib
combination against \textbf{5.4} for chemotherapy \cite{qspmetpancre}, and
RASolute 302 reported \textbf{13.2} against \textbf{6.6 months} in the RAS G12
population nine months later \cite{rasolute302}. That proximity is chronology,
not validation: 1000 simulated against 241 enrolled, a combination against a
single agent, KRAS G12C against a primarily G12D and G12V population.
